\section{Scattering parameters: measuring a network with waves}
\label{sec:sparams}

\noindent\emph{Status: draft -- definitions, redundancy relations and reference conventions are stated and derived;
the numerical verification of the renormalisation formula is done by hand in Section~\ref{sec:sp-refz} and the
power argument of Section~\ref{sec:sp-powerflow} is order-of-magnitude only.}

The previous chapter established that a wire carrying a wave has two travelling components, and that only when
one of them is absent does the ratio of voltage to current at a point equal the characteristic impedance. Every
measurement of a microwave network is a measurement of those travelling waves, and the parameters that describe
the network must therefore be defined on the waves rather than on the terminal currents and voltages. That is
the whole content of scattering parameters, and the reason they are not merely a convenient notation but the
only self-consistent description available once the ports are far enough apart for the phase to matter. This
chapter builds the definition from the power waves of \eqref{eq:sp-ab}, derives the properties that a
measurement can be checked against (\S\ref{sec:sp-properties}), and is explicit about the two conventions that
are usually left implicit and are the source of most of the disagreements in the literature: the reference
impedance (\S\ref{sec:sp-refz}) and the reference plane (\S\ref{sec:sp-refplane}). It closes by asking what
physically moves when power moves (\S\ref{sec:sp-powerflow}), because the answer is surprising enough to be worth
one page and it corrects a picture many of us carry from circuit theory.

The parameters defined here are what the solver of Section~\ref{sec:impl} computes and what the comparisons of
Section~\ref{sec:res-scatter} plot. The one place where the network is not a two-port but a scatterer -- the
radar-cross-section comparison -- reports its answer as a scattering amplitude, which is the same object with
the reference impedance replaced by the impedance of free space.

\subsection{Why impedance and admittance parameters fail at microwave frequencies}
\label{sec:sp-zy}

Impedance parameters are defined by open-circuit measurements. For a two-port,
\begin{equation}
  \begin{pmatrix} V_1 \\ V_2 \end{pmatrix}
  =
  \begin{pmatrix} Z_{11} & Z_{12} \\ Z_{21} & Z_{22} \end{pmatrix}
  \begin{pmatrix} I_1 \\ I_2 \end{pmatrix},
  \qquad
  Z_{ij} = \left.\frac{V_i}{I_j}\right|_{I_k=0,\ k\neq j} ,
  \label{eq:sp-zdef}
\end{equation}
so $Z_{11}$ is the input impedance with port~2 open-circuited. At a millihertz this is easy: switch off the
current, measure the voltage. At a gigahertz the ``open circuit'' is a pair of conductors a fraction of a
wavelength long, and it is an open circuit only if nothing leaves through it -- but the capacitance between the
two conductors of the open end is a real admittance that grows with frequency, and it draws a measurable current.
The condition $I_2=0$ cannot be imposed, so $Z_{11}$ cannot be measured. The same objection applies to
admittance parameters, which need short circuits: a short at the end of a length of line is a short only if its
inductance is negligible, and a $1$ mm wire has about $1$ nH, which is $6.3\ \Ohm$ at 1 GHz and no longer a
short.

There is a second and more fundamental failure, independent of how carefully the terminations are built. The
impedance matrix of an ideal thru connection -- two ports joined by a perfect conductor -- is
\begin{equation}
  Z_{\mathrm{thru}} = \begin{pmatrix} 0 & 0 \\ 0 & 0 \end{pmatrix},
  \label{eq:sp-zthru}
\end{equation}
because the voltage across a perfect short is zero whatever current flows through it. The matrix carries no
information about the thru at all, and the corresponding admittance matrix is singular. A description that
assigns the zero matrix to one of the most useful networks in the subject, and an infinite matrix to the ideal
open, is a description whose numerical values are the wrong quantities to store. What one wants is a parameter
set that is bounded, that exists for every passive network, and that degrades gracefully: a thru should be a
harmless identity, a matched load should be a zero. Scattering parameters are exactly that, and
\eqref{eq:sp-zthru} is the strongest single argument for them.

\subsection{Power waves, and why the normalisation contains a square root}
\label{sec:sp-power}

Let port $i$ have a reference impedance $Z_{0i}$, in general complex, and let $V_i$ and $I_i$ be the phasor
voltage and current flowing into the network at that port. Define the incident and reflected power waves
\begin{equation}
  a_i = \frac{V_i + Z_{0i} I_i}{2\sqrt{\Re Z_{0i}}},
  \qquad
  b_i = \frac{V_i - Z_{0i} I_i}{2\sqrt{\Re Z_{0i}}} .
  \label{eq:sp-ab}
\end{equation}
The definition is invertible in two lines,
\begin{equation}
  V_i = \sqrt{\Re Z_{0i}}\,(a_i + b_i),
  \qquad
  I_i = \frac{a_i - b_i}{\sqrt{\Re Z_{0i}}} ,
  \label{eq:sp-vibreak}
\end{equation}
and it is worth checking that \eqref{eq:sp-vibreak} returns \eqref{eq:sp-ab} identically, since that is the only
algebra in this chapter that is easy to get wrong. For a real reference impedance, $Z_{0i}=Z_0\in\R$, the
definition collapses to the familiar
\begin{equation}
  a_i = \frac{V_i + Z_0 I_i}{2\sqrt{Z_0}},
  \qquad
  b_i = \frac{V_i - Z_0 I_i}{2\sqrt{Z_0}} ,
  \label{eq:sp-ab-real}
\end{equation}
which is the form used throughout this report: every port of every network here is referred to
$50\ \Ohm$ exactly, and the general complex form is quoted only because the renormalisation formula of
\S\ref{sec:sp-refz} is easiest to state with it in view.

Now compute the power entering the port. With the peak-phasor convention used throughout this report -- the same
convention that produced $P=\tfrac12\Re\{VI^*\}$ in Section~\ref{sec:ms-lossless} -- the complex power is
\begin{equation}
  \tfrac12 V_i I_i^{*}
  = \tfrac12\left[\sqrt{\Re Z_{0i}}(a_i+b_i)\right]
            \left[\frac{a_i^*-b_i^*}{\sqrt{\Re Z_{0i}}}\right]
  = \tfrac12\left(|a_i|^2 - |b_i|^2 + b_i a_i^* - a_i b_i^*\right),
  \label{eq:sp-power-general}
\end{equation}
and the last two terms are purely imaginary, $b_ia_i^*-a_ib_i^* = -2j\,\Im\{a_ib_i^*\}$. Taking the real part,
\begin{equation}
  P_i = \Re\left\{\tfrac12 V_i I_i^*\right\} = \tfrac12\left(|a_i|^2 - |b_i|^2\right) .
  \label{eq:sp-power}
\end{equation}
So $|a_i|^2/2$ is the power carried into the port by the incident wave and $|b_i|^2/2$ the power carried out by
the reflected wave, each measured in watts when the phasors are peak values. This is the statement that makes
scattering parameters physical rather than algebraic, and the factor of two in \eqref{eq:sp-power} is the single
most common place a power budget ``loses'' half. With rms phasors the same equation reads $P_i=|a_i|^2-|b_i|^2$
and \eqref{eq:sp-ab} gains no factor, because the $\sqrt{2}$ is absorbed into the definition of the phasor. The
two conventions differ by exactly a factor of two in $|a|^2$, and a chain of reasoning that mixes them will
produce a network that appears to lose 3 dB or to gain 3 dB with no mechanism.

The modulus signs in \eqref{eq:sp-power} deserve a remark. Because the imaginary part of \eqref{eq:sp-power-general}
need not vanish, $|a|^2-|b|^2$ is not in general the complex power; it is its real part, and separately
$|a|^2+|b|^2$ is related to the reactive part. The reason the cross terms cancel for the real part is that a
wave travelling on a line of real characteristic impedance has its voltage and current in phase, so it carries
real power and no reactive power. A wave on a line of complex $Z_0$ carries reactive power even when it travels
in one direction only, which is why the general definition \eqref{eq:sp-ab} carries the $\sqrt{\Re Z_{0i}}$ rather
than $\sqrt{Z_{0i}}$: the normalisation is chosen so that the power identity \eqref{eq:sp-power} survives a
complex reference impedance. That is the reason for the awkward-looking square root of a real part.

\subsection{The definition of $S$, and what a matched termination really is}
\label{sec:sp-definition}

With the waves defined, the scattering matrix of an $n$-port is defined element by element by
\begin{equation}
  S_{ij} = \left.\frac{b_i}{a_j}\right|_{a_k = 0,\ k\neq j} ,
  \qquad\text{i.e.}\qquad
  \begin{pmatrix} b_1 \\ \vdots \\ b_n \end{pmatrix}
  =
  \begin{pmatrix} S_{11} & \cdots & S_{1n} \\ \vdots & & \vdots \\ S_{n1} & \cdots & S_{nn} \end{pmatrix}
  \begin{pmatrix} a_1 \\ \vdots \\ a_n \end{pmatrix}.
  \label{eq:sp-def}
\end{equation}
Read the condition $a_k=0$ carefully, because the whole physical content of the definition is in it. $a_k=0$
means there is no incident wave at port $k$, and the only way to guarantee that is to terminate port $k$ in a
load that emits nothing back: an absorber. Such a termination is called \emph{matched}, and the word is a
statement about the termination, not about the network. A matched load is one that absorbs completely; it has
$b=0$ for every $a$, which by \eqref{eq:sp-power} means it converts every incident watt into heat (or into
whatever else) and reflects nothing. Then and only then does the ratio $b_i/a_j$ isolate one column of $S$.

The practical consequence is that measuring $S_{21}$ requires terminating port~2 well, and the quality of the
whole measurement is the quality of that termination. If the termination reflects $\Gamma_T$, then the wave
$b_2$ returning from it re-enters the network and produces a further contribution at port~1, so the measured
ratio is not $S_{21}$ but
\begin{equation}
  \left.\frac{b_1}{a_1}\right|_{\text{bad termination}}
  = S_{11} + \frac{S_{12}S_{21}\Gamma_T}{1 - S_{22}\Gamma_T} ,
  \label{eq:sp-termination-error}
\end{equation}
a geometric series in the round-trip product $S_{22}\Gamma_T$. The second term is the same multiple-reflection
sum that appears as the de-embedding correction of Section~\ref{sec:ms-deembed}, seen from the other end of the
problem: there the series had to be removed from a measurement of a line, here it has to be suppressed by a good
load. With $|S_{22}\Gamma_T|=0.1$ the error in $S_{11}$ can be a tenth of $S_{21}S_{12}$, which for a
well-matched through structure is of the same order as the quantity being measured.

For a two-port the four elements have direct one-line readings. $S_{11}$ is the input reflection coefficient with
the output matched: $S_{11}=\Gamma_{\mathrm{in}}$ at $Z_0$ reference. $S_{21}$ is the forward transmission
coefficient, the ratio of the outgoing wave at port~2 to the incident wave at port~1. $S_{12}$ is the reverse
transmission, and $S_{22}$ the output reflection. Each is a complex number with a magnitude and a phase, and each
therefore carries two pieces of information; a two-port is eight real numbers, not four. The magnitudes are what
gets plotted as insertion and return loss; the phases are what carries $\beta$, and Section~\ref{sec:ms-estimators}
is entirely about extracting a propagation constant from the phase of $S_{21}$.

\subsection{$S_{11}$ as an input impedance, and the link to transfer matrices}
\label{sec:sp-zin}

Terminate port~2 in $Z_0$, so that $a_2=0$, and look into port~1. From \eqref{eq:sp-vibreak} the input
impedance is $Z_{\mathrm{in}}=V_1/I_1=Z_0(a_1+b_1)/(a_1-b_1)=Z_0(1+S_{11})/(1-S_{11})$, and inverting gives
\begin{equation}
  S_{11} = \frac{Z_{\mathrm{in}}-Z_0}{Z_{\mathrm{in}}+Z_0} = \Gamma_{\mathrm{in}} ,
  \label{eq:sp-s11-gamma}
\end{equation}
which is \eqref{eq:tline-gamma} again. This is a useful thing to have established: $S_{11}$ and $\Gamma$ are not
analogous quantities, they are the same quantity, and the only difference is that $\Gamma$ is defined on a line
where $Z_0$ is a property of the line, while $S_{11}$ is defined on a port where $Z_0$ is a convention. The
corollary is the one that matters for measurement: everything Section~\ref{sec:smith} says about $\Gamma$ at a
reference plane is true of $S_{11}$, including the fact that $|S_{11}|\le1$ for a passive one-port.

For a two-port it is often convenient to switch to the transfer (ABCD) matrix, which chains: cascading two
two-ports multiplies their ABCD matrices. The conversion from ABCD to $S$ at equal real reference impedances is
\begin{equation}
  \begin{aligned}
    S_{11} &= \frac{A + B/Z_0 - C Z_0 - D}{\Delta}, &
    S_{12} &= \frac{2(AD-BC)}{\Delta}, \\
    S_{21} &= \frac{2}{\Delta}, &
    S_{22} &= \frac{-A + B/Z_0 - C Z_0 + D}{\Delta},
  \end{aligned}
  \qquad
  \Delta = A + \frac{B}{Z_0} + C Z_0 + D ,
  \label{eq:sp-abcd-s}
\end{equation}
and it is worth testing \eqref{eq:sp-abcd-s} against the two networks of \S\ref{sec:sp-zy}. The ideal thru has
$A=D=1$, $B=C=0$, so $\Delta=2$ and $S_{11}=S_{22}=0$, $S_{21}=S_{12}=1$: an identity, as required, and in
sharp contrast with \eqref{eq:sp-zthru}. A series impedance $Z$ has $A=D=1$, $B=Z$, $C=0$, so
$\Delta=2+Z/Z_0$ and
\begin{equation}
  S_{11} = \frac{\Delta-2}{\Delta} = \frac{Z}{2Z_0+Z},
         \qquad\text{so}\qquad
         S_{11} = \frac{Z/Z_0}{2+Z/Z_0} = \frac{Z}{Z+2Z_0} ,
  \qquad
  S_{21} = \frac{2Z_0}{Z+2Z_0} ,
  \label{eq:sp-series-z}
\end{equation}
so with $Z=Z_0$ the series element gives $S_{11}=\tfrac13$ and $S_{21}=\tfrac23$, and
$|S_{11}|^2+|S_{21}|^2=\tfrac19+\tfrac49=\tfrac59<1$ -- the missing $\tfrac49$ having gone into the series
resistor, as it must. Both of these are rederived below in \S\ref{sec:sp-refz}, where they serve as the
numerical test of the renormalisation formula.

\subsection{Reciprocity, losslessness and passivity}
\label{sec:sp-properties}

Three properties of a scattering matrix can be checked from its elements alone, and because they are necessary
conditions they are what a computed or measured matrix should be validated against before any further use.

\emph{Reciprocity.} If the network is built from isotropic, linear, passive materials with no gyrotropic and no
active constituents, then the network is reciprocal and
\begin{equation}
  S = S^{\mathsf{T}} ,
  \qquad\text{i.e.}\qquad S_{ij} = S_{ji}\ \ \forall i,j .
  \label{eq:sp-reciprocity}
\end{equation}
This halves the number of independent measurements: $S_{12}$ need not be measured if $S_{21}$ is, and a measured
matrix with $|S_{21}-S_{12}|$ large is either non-reciprocal or wrong. The reason it holds is that the underlying
differential operator is self-adjoint for such materials, which is the same self-adjointness that gives the
eigenvalue problem of Section~\ref{sec:eigen} its real spectrum.

\emph{Losslessness.} If the network contains no dissipative material and no radiating aperture, then power is
conserved and \eqref{eq:sp-power} applied to all ports at once gives
\begin{equation}
  \sum_{k}|b_k|^2 = \sum_{k}|a_k|^2
  \quad\Longrightarrow\quad
  S^{\mathsf{H}} S = I ,
  \label{eq:sp-lossless}
\end{equation}
so the scattering matrix is unitary. The columns of $S$ are then orthonormal: each has unit norm (a lossless
two-port cannot absorb, so whatever enters one port leaves through the two), and distinct columns are orthogonal
(a lossless network fed at port~1 cannot also deliver to the matched load at port~2, at any frequency). Unitarity
is a strong constraint, and a computed lossless structure whose $S$ deviates from it by more than the discretisation
error is a structure that is losing power to the grid.

\emph{Passivity.} If the network can dissipate but cannot generate, then for every set of incident waves
$|b|^2\le|a|^2$, which is to say
\begin{equation}
  S^{\mathsf{H}} S \le I ,
  \label{eq:sp-passivity}
\end{equation}
meaning that $I - S^{\mathsf{H}}S$ is positive semi-definite. Equivalently, every singular value of $S$ is at
most one, so no incident excitation produces a reflected or transmitted wave with more power than it supplied.

A warning is needed here, and it is the most important sentence in this subsection. \emph{Passivity is a
necessary condition and not a sufficient one.} A solver can produce a scattering matrix that satisfies
\eqref{eq:sp-passivity} exactly and is still wrong: the matrix may be too lossy (a scheme with numerical
dissipation of the right sign satisfies \eqref{eq:sp-passivity} comfortably while smearing every resonance), it
may be passive and non-causal, it may be passive and violate reciprocity, and it may be passive while getting the
phase of $S_{21}$ badly wrong -- passivity constrains magnitudes through \eqref{eq:sp-passivity}, and constrains
phases only through their coupling to magnitudes. Passivity is the one property that a badly converged
calculation is most likely to preserve, because damping is the direction numerical error usually errs in. It is
therefore a sanity check that rules out one class of failure and says nothing about the rest; convergence and
independent ground truth are what establish correctness, and the comparisons of Section~\ref{sec:res-scatter}
are built on the latter.

\subsection{The reference impedance is part of the number}
\label{sec:sp-refz}

An $S$-parameter without a stated reference impedance is not a number. The definition \eqref{eq:sp-def} divides
by $\sqrt{\Re Z_{0i}}$, and the condition $a_k=0$ is a statement about a termination of impedance $Z_{0k}$, so
changing $Z_{0k}$ changes every element of the matrix. The reason this is easy to forget is that a great deal of
software reports $S$ referred to $50\ \Ohm$ without saying so, and then a user compares numbers with a paper that
used $75\ \Ohm$ and finds a disagreement that is not an error.

The reason a real reference impedance is the usual choice is worth stating precisely. A real $Z_0$ cannot absorb
a reactance: an incident wave on a line of real $Z_0$ has voltage and current in phase, so the load that reflects
nothing is a pure resistance. Reference a port to a complex $Z_0=R_0+jX_0$ and the matched termination must be
that complex impedance, which is a particular circuit that happens to sit at one frequency only; sweep the
frequency and the ``matched'' load is no longer matched. A real reference impedance is therefore the only choice
that gives a frequency-independent meaning to the word matched, and that is why $50\ \Ohm$ is used even for
structures whose natural impedance is nothing of the kind.

The renormalisation formula is worth having explicitly. Let $\rho=(Z_0'-Z_0)/(Z_0'+Z_0)$ be the reflection
coefficient of the new reference impedance in the old one, and let $S$ be referred to $Z_0$, $S'$ to $Z_0'$, all
ports changed at once and both real. Then
\begin{equation}
  S' = (S - \rho I)(I - \rho S)^{-1} ,
  \label{eq:sp-renorm}
\end{equation}
and for a two-port, writing it out,
\begin{equation}
  \begin{aligned}
    S'_{11} &= \frac{(S_{11}-\rho)(1-\rho S_{22}) + \rho S_{12}S_{21}}{\Delta}, &
    S'_{12} &= \frac{S_{12}(1-\rho^2)}{\Delta}, \\
    S'_{21} &= \frac{S_{21}(1-\rho^2)}{\Delta}, &
    S'_{22} &= \frac{(S_{22}-\rho)(1-\rho S_{11}) + \rho S_{12}S_{21}}{\Delta},
  \end{aligned}
  \label{eq:sp-renorm-2p}
\end{equation}
where
\begin{equation}
  \Delta = (1-\rho S_{11})(1-\rho S_{22}) - \rho^2 S_{12}S_{21} .
  \label{eq:sp-renorm-delta}
\end{equation}
For a one-port \eqref{eq:sp-renorm} collapses to the M\"obius form
\begin{equation}
  \Gamma' = \frac{\Gamma - \rho}{1 - \rho\Gamma} ,
  \label{eq:sp-renorm-1p}
\end{equation}
which is the same bilinear map that will be drawn as the Smith chart in Section~\ref{sec:smith}, and it is a good
check that \eqref{eq:sp-renorm-2p} reduces to it when the ports are decoupled.

Two numerical tests fix the formula. First, take the series impedance of \eqref{eq:sp-series-z} with $Z=Z_0$:
$S_{11}=\tfrac13$, $S_{21}=\tfrac23$, $S_{12}=\tfrac23$, $S_{22}=\tfrac13$. Renormalise from $Z_0$ to $Z_0'=2Z_0$,
so $\rho=\tfrac13$. Then $\Delta=(1-\tfrac19)^2-(\tfrac19)(\tfrac49)=\tfrac{64}{81}-\tfrac{4}{81}=\tfrac{60}{81}$,
and
\begin{equation}
  S'_{11} = \frac{\left(\tfrac13-\tfrac13\right)\left(1-\tfrac19\right)+\tfrac13\cdot\tfrac49}{\tfrac{60}{81}}
          = \frac{\tfrac{4}{27}}{\tfrac{60}{81}} = \frac15 = 0.2 ,
  \qquad
  S'_{21} = \frac{\tfrac23\left(1-\tfrac19\right)}{\tfrac{60}{81}}
          = \frac{\tfrac{16}{27}}{\tfrac{60}{81}} = \frac45 = 0.8 ,
  \label{eq:sp-renorm-check}
\end{equation}
and the direct calculation agrees: referred to $2Z_0$ the series element $Z=Z_0$ sees $4Z_0$ of total
impedance, giving $S'_{11}=Z/(Z+4Z_0)=\tfrac15$ and $S'_{21}=4Z_0/(Z+4Z_0)=\tfrac45$. Second, take the thru,
$S_{11}=S_{22}=0$, $S_{21}=S_{12}=1$, which must be invariant under any renormalisation because a perfect
connection has no impedance to be renormalised: \eqref{eq:sp-renorm-2p} gives $\Delta=1-\rho^2$, then
$S'_{11}=(0-\rho)(1-0)/(1-\rho^2)+\rho(1)(1)/(1-\rho^2)=(-\rho+\rho)/(1-\rho^2)=0$ and
$S'_{21}=1\cdot(1-\rho^2)/(1-\rho^2)=1$, as required. Both checks pass, and the formula is used in this form wherever
this report compares numbers referred to different impedances.

\subsection{Reference planes, and what a length of line does to $S$}
\label{sec:sp-refplane}

The second convention is where the ports are. Moving a reference plane outwards by a length $L$ of lossless line
of propagation constant $\beta$ does not change the network; it changes only the phase with which the waves are
referred to it. The incident wave at the moved plane has travelled an extra $L$ to reach the old plane, so it
picks up $e^{-j\beta L}$; the reflected wave has travelled $L$ out and $L$ back, so it picks up $e^{-2j\beta L}$.
For a two-port with port~1 moved by $L_1$ and port~2 by $L_2$,
\begin{equation}
  S'_{11} = S_{11}e^{-2j\beta L_1}, \quad
  S'_{22} = S_{22}e^{-2j\beta L_2}, \quad
  S'_{21} = S_{21}e^{-j\beta(L_1+L_2)}, \quad
  S'_{12} = S_{12}e^{-j\beta(L_1+L_2)} .
  \label{eq:sp-refplane}
\end{equation}
The often-quoted pair $S_{11}\to S_{11}e^{-2j\beta L}$ and $S_{21}\to S_{21}e^{-j\beta L}$ is the special case
$L_1=L$, $L_2=0$: it is what happens when the same physical length is added at one port only, and quoting it
without that qualification is how a factor-of-two error enters a de-embedding. The general form
\eqref{eq:sp-refplane} is the one this report uses, and it is exactly the per-port shift $d$ of
\eqref{eq:ms-deembed} applied to a scattering matrix rather than to a propagation factor.

Two consequences are worth naming. First, $|S_{21}|$ is unchanged by a length of lossless line, so attenuation
cannot be mistaken for de-embedding error and vice versa: a measured $|S_{21}|$ that falls with length is
dissipation, while a measured $\angle S_{21}$ that falls with length is propagation. Second, the factor of two in
the reflection terms is the same factor of two as in the standing-wave period of \eqref{eq:tline-ripple-period}
and in the Smith-chart rotation of Section~\ref{sec:smith-rotation}. All three say one thing: a reflection
involves a round trip. A reference-plane shift of $\lambda_g/4$ therefore rotates $S_{11}$ by $\pi$ and
$S_{21}$ by $\pi/2$, which is why a de-embedding carried out at one frequency and assumed valid at a nearby one
fails fastest in the reflection coefficient.

\subsection{What actually carries the power}
\label{sec:sp-powerflow}

Everything above treats $|a|^2$ as a power and never asks what is doing the carrying. The answer is not the
electrons, and the orders of magnitude make the point so emphatically that the picture is worth correcting once.

A current density $J$ in a conductor with $n$ free electrons per unit volume and charge $-e$ corresponds to a
drift velocity
\begin{equation}
  v_d = \frac{J}{n e} .
  \label{eq:sp-drift}
\end{equation}
For copper, $n=8.5\times10^{28}\ \mathrm{m}^{-3}$ and $e=1.602\times10^{-19}$ C. At a current density of
$1$ A/mm$^2$, which is $10^{6}$ A/m$^2$ and a perfectly ordinary value for a PCB trace, \eqref{eq:sp-drift} gives
\begin{equation}
  v_d = \frac{10^{6}}{8.5\times10^{28}\times1.602\times10^{-19}}
      \approx 7.3\times10^{-5}\ \mathrm{m/s},
  \label{eq:sp-drift-num}
\end{equation}
that is, $0.073$ mm/s: the electrons move less than a tenth of a millimetre per second. A typical signal of a few
milliamperes in a $0.1$ mm$^2$ trace is an order of magnitude less again, some $7\times10^{-7}$ m/s, and at that
rate an electron takes about a fortnight to traverse a $1$ m cable. The signal, meanwhile, travels at the phase
velocity of the line, $2\times10^{8}$ m/s on a typical substrate. The ratio is
$2\times10^{8}/7.3\times10^{-5}\approx 3\times10^{12}$: the disturbance moves thirteen orders of magnitude faster
than the charges that constitute it.

The reconciliation is that the electrons are not the message. A signal is a rearrangement of charge density and
field, and the rearrangement propagates because each electron nudges its neighbour through the Coulomb field,
which propagates at the speed of light in the surrounding dielectric. The drift velocity is a statement about how
far the carriers get; the signal velocity is a statement about how fast the pattern of pushes propagates, and the
two are related the way the speed of a wave on a rope is related to how far the rope's material moves.

So what carries the power is the field, and the quantitative statement is Poynting's theorem: the power flowing
through a surface is
\begin{equation}
  P = \oint \left(\mathbf{E}\times\mathbf{H}\right)\cdot\dd\mathbf{A} .
  \label{eq:sp-poynting}
\end{equation}
For a TEM mode both $\mathbf{E}$ and $\mathbf{H}$ are purely transverse, so $\mathbf{E}\times\mathbf{H}$ points
along the direction of propagation and its integral over the cross-section is exactly
$\tfrac12\Re\{VI^*\}$ -- the same quantity that appeared as $|a|^2/2-|b|^2/2$ in \eqref{eq:sp-power}. That identity
is not a coincidence: it is the reason the scattering parameters of a TEM line can be defined at all, and it is
why $|a|^2$ deserves to be called a power. It also separates the energy storage cleanly. In a TEM mode the
electric and magnetic energy densities are equal at every point,
$\tfrac12\eps|\mathbf{E}|^2=\tfrac12\mu|\mathbf{H}|^2$, so each holds half of the total stored energy and the
wave carries the two halves along together: neither can be described as the ``cause'' of the other, and a
description that tracks only the voltage is tracking only half the energy. For the quasi-TEM mode of a microstrip
this equality is approximate, which is the field-theoretic statement of the same caveat as
Section~\ref{sec:tline-limits}: a longitudinal field component tilts $\mathbf{E}\times\mathbf{H}$ slightly and a
little power flows in the transverse direction, near the conductor edges, where the mode is least TEM-like.

The practical moral for the measurements of this report is that a probe measures the field at the point it is
placed, not the current in the conductor beneath it. Section~\ref{sec:ms-lossless} already had to take the
complex power as $\tfrac12\Re\{VI^*\}$ on exactly this basis, and the factor of two between peak and rms phasors
-- one half of the total energy in the electric field, one half in the magnetic -- is the same factor of two that
has to be watched in \eqref{eq:sp-power}.
